\section{Is the robot at the edge of the frame?}
We saw in the previous chapter that the boolean functions {\idt tn()}, {\idt tne()}, {\idt te()}, {\idt tse()}, {\idt ts()},
 {\idt tsw()}, {\idt tw()}, {\idt tnw()} can be used in ``if''
 statements to find out the orientation of the \robot. Another boolean function,
 {\idt tb()} (\underline{t}est \underline{b}order), makes it possible to test whether
 the \robot\ is at the edge (the border) of the frame. More precisely {\idt tb()} is {\mcl True}
 if the supposed execution of the command {\idt fd()} would take
 the \robot\ out of the frame in which it moves.  The boolean expression
 {\idt tb()} is {\mcl False} if the supposed execution
 of the command {\idt fd()} moves the \robot\ only inside the frame in which it moves.
For example, we can
 define a command {\idt fdWithoutLeaving()} which orders the \robot\ to execute an {\idt fd()} if this does not
 take it out of the frame:

\input python/FdWithoutLeavingTheFrame.tex

{\sloppy\noindent The empty parentheses of {\mcl not}\ {\idt tb()} are indispensable.
If we forget them and write {\mcl not}\ {\idt tb}, \python\ reports no error,
but the program does not do what we want: {\idt tb} without parentheses denotes the
function itself, and not the value computed by calling it. Now \python\ considers
that a function, like any value that is not zero or empty, counts as {\mcl True}.
The expression {\mcl not}\ {\idt tb} is therefore always {\mcl False} and the \robot\
never moves forward, whether it is at the edge of the frame or not. This is an error that is hard to find
since nothing reports it!\par}

%\pagebreak
\section{``while'' loop}

 A ``while'' loop makes it possible to repeat the execution of a sequence of in\-struc\-tions zero
or more times, testing at each iteration whether to stop or to continue. The English word
``while'' means here `as long as'. The ``while'' loop makes it possible to
repeat the execution of a sequence of instructions as long as the value of a boolean expression
is {\mcl True}. For example, the following program draws a horizontal line across the middle
of the screen by orienting the \robot\ towards the east and repeating {\idt fd()} as long as the \robot\ does not leave
the frame:

 \input python/HorizontalLine.tex

The result of the execution of this program is shown reduced in the figure below,
the frames in which the \robot\ moves and in which error messages are displayed being dotted:

 \vspace{0.4cm}

 \includevspacePostScript[4.4cm]{TraitHorizontal}{0.95}{0.95}

 \vspace{0.4cm}

\noindent  A ``while'' loop always has the following form:

\index{While loop@``while'' loop}\index{Boolean expression}
\begin{center}
\begin{tabular}{@{}|p{12.5cm}|@{}} \hline
{\em ``while'' loop:}\\ \hline \hline
\begin{minipage}[t]{12.5cm}
{\medskip
\mbox{\blc\ \ \ {\mcl while}\hspace{0.3em}\blanc\hspace{0.3em} {\small\rm Boolean expression} \hspace{0.3em}\blanc\hspace{0.3em}{\idt :}}\par
\mbox{\blc\ \ \ \ \ \ \ \small\rm  \ldots\ Repeated instructions (indented) \ldots}\par
\medskip}
\end{minipage}\\ \hline
\end{tabular}
\end{center}

\noindent The sequence of instructions which is repeated as long as the value of
the {Boolean expression} is {\mcl True} is called the {\em body of
the loop}\index{Loop body}. As for the ``for'' loop, the {\mcl while} line
ends with a colon {\idt :} and the instructions of the body of
the loop are indented (shifted to the right), the loop ending at the
first instruction which is no longer indented.
If the {Boolean expression} is {\mcl False} when the execution of the
``while'' loop begins then this execution ends immediately without the body of the loop being
executed. If the {Boolean expression} is always {\mcl True} at each iteration
then the execution of the ``while'' loop never ends!
One can still stop the computer by switching off the power. More re\-com\-mend\-able methods are explained
below.

\underline{A serious error} is to write a loop body which does not
change the value of the Boolean expression. For example if we write
{\mcl while}\ {\mcl not}\ {\idt tb():\ qtr()}, the \robot\
turns on the spot without ever getting closer to the edge, and the computer loops forever.

We preferably use ``for'' loops (whose execution always terminates), reserving
the use of ``while'' loops for the case where the number of it\-er\-a\-tions cannot be computed by an integer expression before
starting to execute the body of the loop.

\begin{exemple}[Frame]{\rm The program below draws a frame a\-round a drawing:

 {\raggedright

 \vspace{0.4cm}

 \includevspacePostScript[4.4cm]{TracerCadre}{0.95}{0.95}

 \vspace{0.4cm}

\input python/DrawFrame.tex
 }
\noindent $\Box$ }\end{exemple}

All the programs we had written until now terminated. This means that
 any execution of the program takes a finite time, which sometimes seems long, but we can be sure that by waiting long
 enough, we will always be able to
 observe a result. If for example we have lifted the pen of the \robot\ and forgotten to lower it
 afterwards, the result is a blank sheet but the execution of the program terminates all
 the same.

 With the ``while'' loop, we can write programs whose execution does not terminate.  This is
 the case for the following program:

 \input python/Loop1.tex

 \noindent The \robot\ turns in the centre of the screen, as long as the computer is running and nobody intervenes.
 It can be stopped by pressing, in the window of the \robot, one of the keys
 \begin{picture}(12,12)(0, 2) %escape key
  \put(0,0){\framebox(12,12)[lb]{}}
  \put(1.5,2){\makebox(0,0)[lb]{\tiny esc}}
  \end{picture}\hspace{0.3ex},
  \begin{picture}(12,12)(0, 2) %key §
  \put(0,0){\framebox(12,12)[lb]{}}
  \put(1.5,2){\makebox(0,0)[lb]{\tiny \S}}
  \end{picture}\hspace{0.3ex} or
  \begin{picture}(12,12)(0, 2) %key #
  \put(0,0){\framebox(12,12)[lb]{}}
  \put(1.5,2){\makebox(0,0)[lb]{\tiny \#}}
  \end{picture}\hspace{0.3ex}.  The execution of the program ends because before
executing a command, the \robot\ checks whether the user has pressed one of
these keys to ask it to stop.

  On the other hand, if we execute the following \python\ program:

 \input python/Loop2.tex

 \noindent the execution loops without stopping even if we press the stop keys
  \begin{picture}(12,12)(0, 2) %escape key
  \put(0,0){\framebox(12,12)[lb]{}}
  \put(1.5,2){\makebox(0,0)[lb]{\tiny esc}}
  \end{picture}\hspace{0.3ex},
  \begin{picture}(12,12)(0, 2) %key §
  \put(0,0){\framebox(12,12)[lb]{}}
  \put(1.5,2){\makebox(0,0)[lb]{\tiny \S}}
  \end{picture}\hspace{0.3ex} or
  \begin{picture}(12,12)(0, 2) %key #
  \put(0,0){\framebox(12,12)[lb]{}}
  \put(1.5,2){\makebox(0,0)[lb]{\tiny \#}}
  \end{picture}\hspace{0.3ex}
since the \robot\ receives
  no command while the computer executes this loop, whose body
  is reduced to {\mcl pass}, an instruction which does nothing (\python\ requires
  at least one instruction in the body of a loop and {\mcl pass} serves
  that purpose).

  To stop the execution of this program, one must go back to the terminal window
  in which it was launched, for example with the command:
  \begin{center}{\idt ./robot.sh book/17\_while/Loop2.py}\end{center}
  \noindent hold down the control key
\begin{picture}(22,12)(0, 2) %Ctrl key
\put(0,0){\framebox(22,12)[lb]{}}
\put(1.5,2){\makebox(0,0)[lb]{\tiny Ctrl}}
\end{picture}\hspace{0.3ex}
and press the key
\begin{picture}(12,12)(0, 2) %key C
\put(0,0){\framebox(12,12)[lb]{}}
\put(1.5,2){\makebox(0,0)[lb]{\tiny C}}
\end{picture}\hspace{0.3ex}. The operating system of the computer then interrupts
the program, whatever it is doing (and the window of the \robot\ closes).

If the program was launched by dragging its icon onto the application \texttt{LeRobotPythonEN.app}, a terminal
window opens, which can be selected with the mouse in order to type Ctrl+C.

%  On the first Macintosh computers, for which this book was originally written, one
%  had to press the interrupt key (marked with a circle)
%  located on the side of the computer (figure \ref{fig:interrupt}) while avoiding
%  the reset key (marked with a triangle), which amounted to switching the
%  Macintosh off and on again.
%
%\begin{figure}[htbp]
%  \includevspacePostScript[2.9cm]{Interrupt}{0.95}{0.95}
%  \caption{Reset and interrupt keys of the first Macintosh computers}\label{fig:interrupt}
%\end{figure}

\clearpage
\needspace{6\baselineskip}\noindent \rule{137mm}{0.1mm}\par\nopagebreak %---------------------------------------------
\begin{exercices}{Write \python\ programs using ``while'' loops to produce the following
drawings:

\tiret  the frieze of figure \ref{fig:FriTri} below;
% (page \pageref{fig:FriTri});

\tiret  the grid of the stopping points of the \robot\ of figure \ref{fig:DgQuGr} (page \pageref{fig:DgQuGr}),
without of course using the command {\idt dg()};

\tiret  the square grid of figure \ref{fig:DgQuGr} (page \pageref{fig:DgQuGr});

\tiret  the lattice of figure \ref{fig:DgQuGr} (page \pageref{fig:DgQuGr}) obtained by superimposing two
square grids, one of which is turned by an eighth turn.

In all cases the drawing must be as large as possible and fit inside the frame set for the
movements of the \robot\ on the computer screen.

\noindent $\Box$}\end{exercices}

\vspace*{-1ex}%
\begin{exercice}{Define the commands {\idt gl()},  {\idt gr()},  {\idt gt()} and  {\idt gb()} by procedures containing
``while'' loops (use the boolean function {\idt valPd()} which is {\mcl True} if and
only if the pen of the \robot\ is down).
\noindent $\Box$}\end{exercice}

\begin{figure}[htbp]
 \includevspacePostScript[4.4cm]{FriseDeTriangles}{0.9}{0.9}
 \caption{Frieze of triangles (on a reduced screen)}\label{fig:FriTri}
\end{figure}
\begin{figure}[htbp]
 \includevspacePostScript[14.6cm]{DgQuadGrillage}{0.95}{0.95}
 \caption{Stopping points, square grid and lattice on a reduced screen}\label{fig:DgQuGr}
\end{figure}

%%%%%%%%%%%%%%%%%%%%%%%%%%%%%%%%%%%%%%%%%%%%%%%%%%%%%%%%%
\vspace*{-1ex}%
\begin{corrige}{
We give one program, the others being in the folder {\idt book/\allowbreak 17\_while}:
\input python/FriezeOfTriangles.tex
\noindent $\Box$}\end{corrige}
\begin{corrige}{
We give the solution for the command {\idt gl}, the others being in the program
{\idt GlGrGtGb} of the folder {\idt book/\allowbreak 17\_while}:
\enlargethispage{2\baselineskip}
\input python/GlGrGtGb.tex
\noindent $\Box$}\end{corrige}

%%%%%%%%%%%%%%%%%%%%%%%%%%%%%%%%%%%%%%%%%%%%%%%%%%%%%%%%%%%%%%%%%%%%%%%%%%%
\noindent \framebox[\textwidth]{
\begin{minipage}{\myparboxwidth}
\setlength{\parindent}{2em}
\small \sf
\vspace{\abovedisplayskip}
\section*{Addenda}

 In other programming languages, such as \pascal, there is another form of
 unbounded repetition, the ``repeat'' loop.
 It has the following form:

 \index{While loop@``while'' loop}
\begin{center}
\begin{minipage}{10cm}
{\raggedright
\mcl repeat\\
\blc\ \ \small\rm  \ldots\ Repeated instructions (separated by `$;$') \ldots\\
\mcl until\hspace{0.3em}\blanc\hspace{0.3em} \small\rm Boolean expression\\
}\end{minipage}
\end{center}

\noindent The body of the loop is repeated ({\idt repeat}) until ({\idt until}) the value of
 the {\small\rm Boolean expression} is equal to {\idt true}. It is therefore always executed
 at least once.
%\vspace{\abovedisplayskip}
%\end{minipage}
%}
%%%%%%%%%%%%%%%%%%%%%%%%%%%%%%%%%%%%%%%%%%%%%%%%%%%%%%%%%%%%%%%%%%%%%%%%%%%%

%%%%%%%%%%%%%%%%%%%%%%%%%%%%%%%%%%%%%%%%%%%%%%%%%%%%%%%%%%%%%%%%%%%%%%%%%%%
%\noindent \framebox[\textwidth]{
%\begin{minipage}{\myparboxwidth}
%\setlength{\parindent}{2em}
%\small \sf
%\vspace{\abovedisplayskip}
\indent The ``repeat'' loop does not exist in \python, but it can always be expressed
with a ``while'' loop. A ``repeat'' loop:
 {
\parindent \parindentprog
\medskip

\pagebreak[3]
\mbox{\blc  \blk\mcl repeat}\par
\mbox{\blc  \blk\blk\small\rm Body of the loop}\par
\mbox{\blc  \blk\mcl until {\small\rm Boolean expression}}\par

\medskip }

\noindent can be written in \python\ with a ``while'' loop as follows:
{
\parindent \parindentprog
\medskip

\mbox{\blc  \blk\small\rm Body of the loop}\par
\mbox{\blc  \blk\mcl while\ not\ \small\rm Boolean expression\idt :}\par
\mbox{\blc  \blk\blk\small\rm Body of the loop}\par

\medskip }

\noindent but the body of the loop must then be written twice. In \python\ one often
prefers to write a ``while'' loop whose test is always
{\mcl True} and which is exited with the instruction {\mcl break}\index{break},
which immediately terminates the execution of the loop:
{
\parindent \parindentprog
\medskip

\mbox{\blc  \blk\mcl while\ True\idt :}\par
\mbox{\blc  \blk\blk\small\rm Body of the loop}\par
\mbox{\blc  \blk\blk\mcl if\ \small\rm Boolean expression\idt :}\par
\mbox{\blc  \blk\blk\blk\mcl break}\par

\medskip }

\vspace{\abovedisplayskip}
\end{minipage}
}

%%%%%%%%%%%%%%%%%%%%%%%%%%%%%%%%%%%%%%%%%%%%%%%%%%%%%%%%%%%%%%%%%%%%%%%%%%%%
%%%%%%%%%%%%%%%%%%%%%%%%%%%%%%%%%%%%%%%%%%%%%%%%%%%%%%%%%%%%%%%%%%%%%%%%%%%
\noindent \framebox[\textwidth]{
\begin{minipage}{\myparboxwidth}
\setlength{\parindent}{2em}
\small \sf
\vspace{\abovedisplayskip}
\noindent  Likewise, a ``while'' loop:

{
\parindent \parindentprog
\medskip

\mbox{\blc  \blk\mcl while\ \small\rm Boolean expression\idt :}\par
\mbox{\blc  \blk\blk\small\rm Body of the loop}\par

\medskip }

\noindent could be rewritten with a ``repeat'' loop as follows:

{
\parindent \parindentprog
\medskip

\mbox{\blc  \blk\mcl if\ \small\rm Boolean expression\ \mcl then}\par
\mbox{\blc  \blk\blk\mcl repeat}\par
\mbox{\blc  \blk\blk\blk\small\rm Body of the loop}\par
\mbox{\blc  \blk\blk\mcl until\ not\ \idt (\small\rm Boolean expression\idt )}\par

\medskip }

The use of one or the other form of iteration depends only on whether the body of the loop
may never be executed (``while'' loop) or must always be executed at least once
(``repeat'' loop). The designers of \python\ decided that the ``while'' loop was enough.

\vspace{\abovedisplayskip}
\end{minipage}
}
%%%%%%%%%%%%%%%%%%%%%%%%%%%%%%%%%%%%%%%%%%%%%%%%%%%%%%%%%%%%%%%%%%%%%%%%%%%%
